\newpage

% =============================================================================
% SENTEZ MALZEMESI  (analizden cikan, tek bolumde soylenemeyen bulgular)
%
% Ayrimin kurali: bu cumleyi TEK BIR DILE bakarak soyleyebiliyor muyum?
% Evetse o dilin bolumune ait. Hayirsa buraya.
%
% -----------------------------------------------------------------------------
% 1. IKAME. Iki kayit ters yone gidiyor ama AYNI ANDA DEGIL.
%      Yolsuzluk ve cete   kirilma 2023-11   -66%  (lider degisimiyle ayni ay)
%      Erdogan ve AKP      kirilma 2024-09  +169%  (on ay sonra)
%    Ikisi de q<0.001, aylik seride, basamak (kayma degil).
%    Yani "yeni lider yeni dil" degil: eski dil liderle birakildi, yeni dil
%    on ay sonra ayri bir konjonkturde basladi. Iki ayri karar.
%
% 2. KONFAUNDA KARSI KANIT. Hareket etmeyenler hareket edenler kadar onemli:
%      Ekonomik talepler   +3%
%      Kadin (anma)        -0%
%    Genel bir uslup kaymasi olsaydi hepsi birlikte oynardi. Yolsuzluk %64
%    duserken ekonomik dilin hic kimildamamasi, degisimin SECICI oldugunu
%    gosteriyor. Bu ancak oynayanlarla oynamayanlar yan yana konunca gorunur.
%
% 3. MEKAN BULGULARI TASIMIYOR. Son tabakalama ile mekan bilesimi iki donemde
%    esitlendi (agirliklar: meclis %25, etkinlik %24, tv %20, miting %16,
%    basin %14, kurultay %1). On uc kayittan on birinde sonuc degismedi:
%      Erdogan ve AKP     +176% -> +177%
%      Yolsuzluk ve cete   -64% ->  -64%
%      Ahlaki hak edis     -77% ->  -78%
%    Iki istisna (Ideolojik kimlikler, Kadin haklari) ve ikisinde de duzeltme
%    etkiyi BUYUTUYOR. Yani mekan bulgu yaratmiyor, gizliyor.
%    Bu tek bolumde soylenirse on uc kez tekrarlanir, burada bir kez.
%
%
% -----------------------------------------------------------------------------
% YAPISAL NOT
% Bu bolumun adi su an "Bulgular" ama ondan onceki bes bolum de bulgu.
% Ad, bolumun ne YAPTIGINI soylemeli: "Diller Arasinda Tekrar Edenler" gibi.
% Yeri de leksikal blok ile semantik blok arasinda kaliyor; bir sentez
% sentezledigi seyden sonra gelmeli, yani 09_erdogan'in ardina.
%
% Asagidaki birinci tablo tamamen "Ekonomik talepler" satirlarindan olusuyor,
% yani tek dilin olay davranisi: 02_ekonomik'e ait, buraya degil.
% =============================================================================

\section{Bulgular}

% Her bulgunun basligi IDDIANIN KENDISI. Basliklar icindekiler listesine
% dusuyor, yani okuyucu sadece icindekilere bakarak argumani gorebiliyor.

\begin{bulgu}{Yolsuzluğa ve çeteye karşı ahlaki çerçeve sönümleniyor. Etnik ve inanç kimlikleri
söylemde yer ediniyor.}

Kul hakkı gasp edilmiş ve hak, hukuk, adalet mücadelesi veren insanlara karşı beşli çete, mafya ve baron
gibi figürlerle doldurulan düşman anlatısı Özgür Özel döneminde terk ediliyor. 
Erdoğan ve AKP kelimelerinin frekansı 2024 Eylül'de 2.7 kat artıyor. Etnik ve inanç kimlikleri
2024 yerel seçimlerinden sonra CHP söyleminde düzenli bir biçimde çağrılmaya başlanıyor. 


\begin{table}[H]
\centering
\caption{Kayıtların iki dönem arasındaki değişimi}
\label{tab:register-change}
\begin{tabular}{@{}l r r r c@{}}
\toprule
 & Kılıçdaroğlu & Özel & Değişim (kat) & \\
\midrule
Etnik ve inanç kimlikleri & 0.6 & 7.6 & +12.5x & \artan \\
Erdoğan ve AKP & 16.5 & 45.5 & +2.8x & \artan \\
Ahlaki çerçeve & 11.1 & 2.5 & -4.3x & \azalan \\
Yolsuzluk ve çete & 14.9 & 5.3 & -2.8x & \azalan \\
\bottomrule
\end{tabular}
\kaynaknot{...}
\end{table}

\end{bulgu}

\begin{bulgu}{Ekonomik talepler ve sınıflar söylemin en kırılgan(?) olabileceği anlara duyarlı. İki kategori CHP
  söyleminde sürekli kesintiye uğruyor.} 
Siyasi bir rejimin sıklıkla deneyimlemediği ve tarafların yeni bir söylem geliştirmek zorunda
kaldığı anlar (dislokasyon), diğer temalara kıyasla sistematik bir şekilde ekonomik talepleri ve sınıfları etkiliyor.
İktidarın etkisinin olduğu dislokasyonlarda ekonomik taleplerin eski seviyelerine dönmesi darbe ve anayasa
referandumuna kıyasla 2 ila 3 ay daha uzun sürüyor. CHP söylemi darbe girişiminden sonraki üç aylık süreçte ekonomik talepleri ve rakibe yönelik dili (örn. diktatör) bırakıyor, buna karşılık kuruluş anlatısını ve ideolojik kimlikleri öne çıkarıyor. 
  % burada en kırılgandan kastım dislokasyon olduğunu söylemdeden anlatmak. yani bu anlar
  % söylemin yeniden kurgulanmak zorunda kalındığı anlar



\begin{table}[htbp]
\centering
\small
\setlength{\tabcolsep}{4pt}
\caption{Log-likelihood keyness of CHP discourse around seven critical events (signifier-independent): words that \textsc{enter} (keyness $+$) vs.\ \textsc{recede} (keyness $-$) across the event, in symmetric $\pm 3$-month windows ($\geq 5$ speeches). Events marked $\dagger$ are treated as dislocations. Each event reactivates a \emph{distinct} register rather than following a single law. Lemma artefacts merged (\textit{üre}$\rightarrow$\textit{üretim}); place-names omitted.}
\label{tab:dislocation-all}
\begin{threeparttable}
\begin{tabular}{@{}>{\raggedright\arraybackslash}p{2.6cm} >{\raggedright\arraybackslash}p{5.3cm} >{\raggedright\arraybackslash}p{5.3cm}@{}}
\toprule
\textbf{Dislokasyon} & \textbf{Enters (reactivated)} & \textbf{Recedes} \\
\midrule
Darbe girişimi†, 2016 & darbe\tnote{a}, OHAL, liyakat, cumhuriyet, laik & dokunulmazlık, üretim\tnote{b}, asgari\_ücret, lira, elektrik, \textbf{diktatör}, turizm, [sosyal] yardım \\
\addlinespace
Referandum†, 2017 & adalet, darbe, savcı, yürüyüş, gayrımeşru, meşru, saray & anayasa, hayır, millet, ulusal, \textbf{tek\_adam\_rejimi}, otorite, Suriye \\
\addlinespace
Genel seçimler, 2018 & \textbf{kriz}, ihale, merkez\_bankası, israf, dolar, borç & esnaf, şeker [fabrika], barışş, kadın, eğitim, emekli, okul, kadın, \textbf{vesayet}, aday \\
\addlinespace
Yerel seçimler, 2019† & adalet, mazbata, hakim, vicdan, iptal & milliyetçi, hizmet, \textbf{özelleştirme}, lira, \textbf{fabrika}, kuyruk, belediye, mahalle, vergi \\
\addlinespace
Genel seçimler, 2023 & \textbf{Erdoğan}, borç, gazeteci, ahlak, dolandırıcı, montaj, kurultay, medya, irade & kul\_hakkı, emekli, genç, beşli\_çete, deprem, kuruş, öğretmen, adalet, okul, vaat kelimeleri\tnote{c} \\
\addlinespace
Yerel seçimler, 2024 & anayasa, birinci, \textbf{Gezi}, \textbf{normalleşme}, müzakere, IMF, müfredat, buğday, kamu & emekli, ittifak, aday, lira, altın, saray, DEM Parti, Erdoğan \\
\addlinespace
İmamoğlu tutuklanması 2025† & İmamoğlu, Saraçhane, sivil\_darbe, cunta, kayyum, boykot, genç, terörsüz, Silivri, rüşvet, diploma & asgari\_ücret, geçim, emekli, maaş, kadın, orman\_yangını, deprem\tnote{d} \\
\bottomrule
\end{tabular}
\begin{tablenotes}\footnotesize
\item[a] Ve FETÖ, cemaat, temmuz ve girişim gibi kelimeler.
\item[b] Ve çiftçi, tarım, fındık ve ürün gibi üretimle ilgili kelimeler.
\item[c] Söz, getireceğim, getireceğiz, yapacağız, göreceksiniz\dots
\item[d] Ve 6 Şubat 2023 depremleriyle ilgili kelimeler: Adıyaman, Hatay, Kahramanmaraş, container housing.
\end{tablenotes}
\end{threeparttable}
\end{table}


\end{bulgu}

\begin{bulgu}{Geçmişin antagonistleri (buraya başlık lazım bir de ne tartışacağım)}

\begin{table}[H]
\centering
\caption{Kayıtların iki dönem arasındaki değişimi}
\label{tab:register-change-akp}
\begin{tabular}{@{}l r r r c@{}}
\toprule
 & Kılıçdaroğlu & Özel & Değişim (kat) & \\
\midrule
Darbe çerçevesi & 10.7 & 16.3 & +1.5x & \artan \\
Cemaat & 5.0 & 1.8 & -2.7x & \azalan \\
Terör & 11.1 & 5.3 & -2.1x & \azalan \\
\bottomrule
\end{tabular}
\kaynaknot{...}
\end{table}

\end{bulgu}

\begin{bulgu}{Emek}

%Emek ve çalışma & 23.1 & 40.9 & +1.8x & \artan \\

\end{bulgu}

% burayı hegemonik analiz olarak yorumlayamam gibime geliyor tartışacağım
%\multicolumn{5}{@{}l}{\textbf{Değişmeyenler}} \\
 % Ekonomik talepler & 129.3 & 133.4 & 1.0x & \duraga \\
  %Kadın hakları & 13.1 & 13.0 & 1.0x & \duraga \\
  %İdeolojik kimlikler & 6.3 & 5.6 & 1.1x & \duraga \\
  %Kurtuluş anlatısı & 5.5 & 4.6 & 1.2x & \duraga \\







